\documentclass[aspectratio=169]{beamer}

% Theme & colors
\usetheme{Madrid}
\usecolortheme{default}
\useinnertheme{circles}

\definecolor{primary}{RGB}{25, 70, 140}
\definecolor{secondary}{RGB}{240, 245, 255}
\definecolor{accent}{RGB}{200, 20, 40}

\setbeamercolor{palette primary}{bg=primary, fg=white}
\setbeamercolor{palette secondary}{bg=primary!80, fg=white}
\setbeamercolor{palette tertiary}{bg=primary!60, fg=white}
\setbeamercolor{palette quaternary}{bg=primary, fg=white}
\setbeamercolor{structure}{fg=primary}
\setbeamercolor{block title}{bg=primary, fg=white}
\setbeamercolor{block body}{bg=secondary}
\setbeamercolor{alerted text}{fg=accent}
\setbeamercolor{frametitle}{bg=primary, fg=white}

% Clean up footer
\setbeamertemplate{navigation symbols}{}
\setbeamertemplate{footline}{%
  \leavevmode%
  \hbox{%
    \begin{beamercolorbox}[wd=.33\paperwidth,ht=2.4ex,dp=1ex,center]{palette primary}%
      \usebeamerfont{author in head/foot}\insertshortauthor
    \end{beamercolorbox}%
    \begin{beamercolorbox}[wd=.34\paperwidth,ht=2.4ex,dp=1ex,center]{palette secondary}%
      \usebeamerfont{title in head/foot}\insertshorttitle
    \end{beamercolorbox}%
    \begin{beamercolorbox}[wd=.33\paperwidth,ht=2.4ex,dp=1ex,right]{palette primary}%
      \usebeamerfont{date in head/foot}\insertframenumber{}/\inserttotalframenumber\hspace*{2ex}
    \end{beamercolorbox}}%
}

% Packages
\usepackage{booktabs}
\usepackage{graphicx}
\usepackage{amsmath, amssymb}
\usepackage{tikz}
\usetikzlibrary{shapes.geometric, arrows.meta, positioning}
\usepackage{xcolor}

% Metadata
\title[Short Title]{Protocolos 4to6 e 6to4}
\subtitle{Por que? Como? E seus impactos para a QoS}
\author[A. Author]{Miguel Moret}
\institute[University]{FCT-UNESP}
\date{\today}

% ─────────────────────────────────────────────
\begin{document}

% ── Title Slide ──────────────────────────────
\begin{frame}
  \titlepage
\end{frame}

% ── Table of Contents ────────────────────────
\begin{frame}{Outline}
  \tableofcontents
\end{frame}

% ═════════════════════════════════════════════
\section{Introduction \& Motivation}
% ═════════════════════════════════════════════

\begin{frame}{Conceito \& Motivação}
  \begin{columns}[T]
    \begin{column}{0.55\textwidth}
      \textbf{Why does this matter?}
      \begin{itemize}
        \item Context and significance of the problem
        \item Existing gap in the literature
        \item Real-world impact or application
      \end{itemize}
      \vspace{0.8em}
      \textbf{Research Question}
      \begin{alertblock}{}
        \centering\textit{``What is the precise question your work addresses?''}
      \end{alertblock}
    \end{column}
    \begin{column}{0.4\textwidth}
      \begin{block}{Key Facts}
        \begin{itemize}
          \item Statistic or figure \#1
          \item Statistic or figure \#2
          \item Statistic or figure \#3
        \end{itemize}
      \end{block}
    \end{column}
  \end{columns}
\end{frame}


% ═════════════════════════════════════════════
\section{Funcionamento do 4to6}
% ═════════════════════════════════════════════

\begin{frame}{Funcionamento do 4to6}
  \begin{columns}[T]
    \begin{column}{0.48\textwidth}
      \begin{block}{Approach A}
        Prior work focused on \ldots\ achieving \ldots\ but limited by \ldots
      \end{block}
      \vspace{0.5em}
      \begin{block}{Approach B}
        Alternative method used \ldots\ with advantage of \ldots\ yet suffering from \ldots
      \end{block}
    \end{column}
    \begin{column}{0.48\textwidth}
      \begin{block}{Gap We Address}
        \begin{itemize}
          \item Limitation 1 of prior work
          \item Limitation 2 of prior work
          \item Our contribution fills this gap
        \end{itemize}
      \end{block}
    \end{column}
  \end{columns}
\end{frame}
\begin{frame}
    
% ═════════════════════════════════════════════
\section{Funcionamento do 6to4}
% ═════════════════════════════════════════════
{Funcionamento do 6to4}
  \begin{columns}[T]
    \begin{column}{0.48\textwidth}
      \begin{block}{Approach A}
        Prior work focused on \ldots\ achieving \ldots\ but limited by \ldots
      \end{block}
      \vspace{0.5em}
      \begin{block}{Approach B}
        Alternative method used \ldots\ with advantage of \ldots\ yet suffering from \ldots
      \end{block}
    \end{column}
    \begin{column}{0.48\textwidth}
      \begin{block}{Gap We Address}
        \begin{itemize}
          \item Limitation 1 of prior work
          \item Limitation 2 of prior work
          \item Our contribution fills this gap
        \end{itemize}
      \end{block}
    \end{column}
  \end{columns}
\end{frame}

% ═════════════════════════════════════════════
\section{Impacto na Qualidade de Serviço (QoS)}
% ═════════════════════════════════════════════

\begin{frame}{Impacto na Qualidade de Serviço (QoS)}
  \begin{columns}[T]
    \begin{column}{0.48\textwidth}
      \begin{block}{Approach A}
        Prior work focused on \ldots\ achieving \ldots\ but limited by \ldots
      \end{block}
      \vspace{0.5em}
      \begin{block}{Approach B}
        Alternative method used \ldots\ with advantage of \ldots\ yet suffering from \ldots
      \end{block}
    \end{column}
    \begin{column}{0.48\textwidth}
      \begin{block}{Gap We Address}
        \begin{itemize}
          \item Limitation 1 of prior work
          \item Limitation 2 of prior work
          \item Our contribution fills this gap
        \end{itemize}
      \end{block}
    \end{column}
  \end{columns}
\end{frame}


\begin{frame}{Future Work}
  \begin{columns}[T]
    \begin{column}{0.32\textwidth}
      \begin{block}{\centering Short-term}
        \centering
        \vspace{0.3em}
        Improve X by incorporating Y data or technique.
        \vspace{0.3em}
      \end{block}
    \end{column}
    \begin{column}{0.32\textwidth}
      \begin{block}{\centering Medium-term}
        \centering
        \vspace{0.3em}
        Extend to domain Z or scale to larger benchmark.
        \vspace{0.3em}
      \end{block}
    \end{column}
    \begin{column}{0.32\textwidth}
      \begin{block}{\centering Long-term}
        \centering
        \vspace{0.3em}
        Open research question that this work motivates.
        \vspace{0.3em}
      \end{block}
    \end{column}
  \end{columns}
  \vspace{1em}
  \centering
  \large\textbf{Thank you!}\\[0.5em]
  \normalsize Questions? \quad \texttt{your.email@institution.edu}
\end{frame}

% ── References ───────────────────────────────
\begin{frame}[allowframebreaks]{References}
  \begin{thebibliography}{9}
    \bibitem{ref1} Author, A. (Year). \textit{Title of Paper}. Journal Name, vol(issue), pp.
    \bibitem{ref2} Author, B. \& Author, C. (Year). \textit{Another Paper Title}. Conference Name.
    \bibitem{ref3} Author, D. (Year). \textit{Book Title}. Publisher.
  \end{thebibliography}
\end{frame}

\end{document}